% Reader question: what differs between a temporary role and a continuous
% person? Claim: a role is a slot that keeps no history of its fillers; a
% person is one identity that carries witnessed outcomes across role changes.
% Evidence: the chapter's worked identity kappa_7, its four role intervals
% (navigator, cartographer, reviewer, lookout) and its eight witnessed
% outcomes o_1..o_8 (Sec. role-person). Must distinguish: slot continuity
% (the reviewer role survives, but is empty of history) from identity
% continuity (kappa_7 carries the witnesses). Grammar: two aligned timelines
% on one grid, rejected: a Venn diagram (identity is not overlap, it is one
% continuous rail) and two portraits (no comparison field, no time axis).
% Five-second test: a reader points at the lower rail and says "one thread,
% four roles, eight witnesses."
\begin{figure}[H]
\centering
\pdfigurehue{pdviolet}%
\begin{tikzpicture}[pd figure,x=1cm,y=1cm]
  \def\x0{3.00}  \def\xe{10.20}
  % ---- ROLE panel: a slot, refilled by unrelated bodies --------------------
  \node[pd title,anchor=west] at (0.00,7.40) {role: one slot, replaceable fillers};
  \node[pd label,anchor=east,align=right,text width=4.1cm] at ({\x0-.20},6.10)
    {reviewer role\\obligation, capability, authority};
  \foreach \x/\body in {4.30/A,6.10/B,7.90/C,9.70/D} {
    \draw[pd rule] (\x,6.55) -- (\x,6.475);
    \draw[pd rule] (\x,5.725) -- (\x,5.65);
    \node[pd state,circle,minimum size=7.5mm,inner sep=0pt] (r\body) at (\x,6.10) {\body};
  }
  \draw[pd hairline,line width=2.4pt] (\x0,6.10) -- (rA.west);
  \draw[pd hairline,line width=2.4pt] (rA.east) -- (rB.west);
  \draw[pd hairline,line width=2.4pt] (rB.east) -- (rC.west);
  \draw[pd hairline,line width=2.4pt] (rC.east) -- (rD.west);
  \draw[pd hairline,line width=2.4pt] (rD.east) -- (\xe,6.10);
  \node[pd note,anchor=west,text width=7.2cm,align=left] at (\x0,5.25)
    {the slot survives; no outcome history moves between fillers};
  % ---- PERSON panel: one identity crossing role boundaries -----------------
  \node[pd title,anchor=west] at (0.00,4.45) {person: one identity, changing roles};
  \node[pd focus label,anchor=east,align=right,text width=3.3cm] at ({\x0-.20},3.10)
    {$\kappa_7$\\non-forgeable identity};
  \draw[pd focus rule,line width=2.0pt] (\x0,3.10) -- (\xe,3.10);
  \foreach \x in {3.00,4.80,6.60,8.40,10.20} \draw[pd guide] (\x,3.00) -- (\x,1.95);
  \foreach \a/\role in {3.90/navigator,5.70/cartographer,7.50/reviewer,9.30/lookout}
    \node[pd label,anchor=north] at (\a,2.80) {\role};
  \foreach \x/\mark in {3.45/$o_1$,4.35/$o_2$,5.25/$o_3$,6.15/$o_4$,7.05/$o_5$,
                        7.95/$o_6$,8.85/$o_7$,9.75/$o_8$} {
    \draw[pd focus rule] (\x,3.25) -- (\x,2.95);
    \node[pd label,anchor=north] at (\x,1.75) {\mark};
  }
  \node[pd focus label,anchor=west,text width=7.2cm,align=left] at (\x0,0.80)
    {witnessed outcomes accumulate on the identity across every role change};
  % ---- one shared time axis --------------------------------------------------
  \draw[pd rule,-{Stealth[length=1.9mm,width=1.2mm]}] (\x0,-0.15) -- (\xe,-0.15);
  \node[pd label,anchor=east] at ({\x0-.15},-0.15) {earlier};
  \node[pd label,anchor=west] at ({\xe+.15},-0.15) {later};
\end{tikzpicture}
\caption{\textbf{A role is a slot; a person is a longitudinal record.} The upper
track holds one reviewer role constant while unrelated bodies fill it in sequence:
the slot has obligations, capability, and authority, but no reputation of its own.
The lower track follows one identity $\kappa_7$ through four role intervals while
witnessed outcomes $o_1\ldots o_8$ accumulate on the same thread. A respawn may
replace the body, and a reassignment may replace the role, without replacing the
accountable person. [internal]}
\label{fig:role-person}
\end{figure}
